%%%%%%%% CONDITIONAL REVISION -- NOT PART OF THE SUBMISSION %%%%%%%%
% Text that Appendix~\ref{sec:bridge} of doc/manuscript.tex would carry IF the
% conjectures H1-H4 below are proven.  The submission keeps the Status (C)
% version; see attic/submission_docs/BRIDGE_REVISION_IF_PROVEN.md (retired 19 Sep 2026) for the assumption ->
% obligation map and for the exact swap procedure.
\section{Cellular automata and quantum formalism}\label{sec:bridge}
\emph{Status (M, derived).}  The mapping below is a derivation rather than an
analogy: the reversible fragment of the readable dynamics is a permutation on the
ontological states, the dissipative steps act as the coarse-graining that defines
measurement, and the Born weights follow from the counting measure on the fibres
of that coarse-graining.  The inputs that carry the weight are stated as explicit
hypotheses (H1-H4 below) and proven in the accompanying derivations; the two
numerical consequences are the equidistant spectrum of
Sect.~\ref{bridge:spectrum} and the no-signalling corollary of
Sect.~\ref{bridge:nosignal}.

\subsection{The ontological basis}\label{bridge:basis}
For local events we follow 't Hooft's cellular-automaton world-model approach
\cite{thooft}.  An ontological state is one complete configuration of the fabric
--- one fixed-length word per cell (Sect.~\ref{subsec:The-cell-structure}) --- and
it does not superpose: the transition rule maps it to exactly one successor.
Following the operator-mechanics recipe of \cite{thooft,Elze2019,rizzo2020perturbing},
the diagonalisation of the step operator yields the Hamiltonian; in what follows
that recipe is \emph{carried out} for this rule set rather than posited.

\textbf{Lemma 1 (orthonormal basis).}  The set of ontological states admits an
orthonormal basis of a finite-dimensional Hilbert space.
\emph{Proof.}  The state space is finite: a bounded number of bits per cell on a
finite lattice (Sect.~\ref{subsec:The-cell-structure}).  The indicator functions
$\ket{A}=e_{A}$ of the configurations $A$ are therefore orthonormal in the
counting inner product $\braket{A}{B}=\delta_{AB}$, and they span the space.
\hfill$\blacksquare$

Finiteness of the fabric is used only here; it is what keeps the operator algebra
below finite-dimensional.  A template is a complex linear combination
\begin{equation}
\ket{\psi}=\sum_{A}\lambda_{A}\ket{A},
\qquad \sum_{A}|\lambda_{A}|^{2}=1 ,
\end{equation}
read as an observer-level bookkeeping device over ontological states, not as an
object in the ontology.

\subsection{The reversible fragment is a permutation}\label{bridge:permutation}
\textbf{H1 (fragment boundary).}  The readable dynamics splits into a fragment
that is a bijection on the ontological states and a fragment that is not: the
local wave propagation and the polarization reconstruction are bijective, while
the sieve, the collapse-and-reissue and the overwrite steps are many-to-one.

\textbf{Theorem 2 (unitary step and Hamiltonian).}  Let $U$ be the one-step map of
the bijective fragment.  Then $U$ is a permutation matrix and hence unitary;
$U=V\Lambda V^{\dagger}$ with $|\lambda_{i}|=1$; and the Schr\"odinger equation
$i\hbar\,\partial_{t}\ket{\psi}=\hat{H}\ket{\psi}$ holds with
$\hat{H}=\tfrac{i\hbar}{T}\ln U$.
\emph{Proof.}  A permutation matrix contains exactly one $1$ in every row and
column, so $U^{\dagger}U=\mathbb{1}$.  Because the state space is finite, $U$ has
finite order $N$ and its eigenvalues are $N$-th roots of unity, $|\lambda_{i}|=1$;
therefore $\ln U$ is skew-Hermitian and $\hat{H}$ is Hermitian and generates the
stated evolution. \hfill$\blacksquare$

The earlier objection --- that the automaton's dissipative steps have no place in
a unitary step --- is answered by H1 itself: those steps are \emph{outside} $U$.
They enter below as the measurement channel, and it is their absence from $U$ that
makes $U$ a permutation rather than merely a contraction.

\subsection{Equidistant spectrum}\label{bridge:spectrum}
\textbf{Corollary 3 (spectrum).}  The energy levels of the reversible fragment are
equally spaced, $\Delta E=2\pi\hbar/(NT)$, and every spectrum is a set of integer
multiples of one fundamental quantum.
\emph{Proof.}  Immediate from Theorem~2: the eigenvalues are $N$-th roots of
unity, so $\hat{H}$ has $N$ equally spaced levels on a circle of length
$2\pi\hbar/T$. \hfill$\blacksquare$

This is the structural reason the model's bosonic excitations carry an
even-integer ladder rather than a continuum (Sect.~\ref{bosons}, where the pair
stack of frequency $2n$ is the superposition of $n$ coincident pairs), and it is a
falsifiable statement about any system built from this fragment: a non-uniform
(hydrogen-like) spectrum cannot be generated by $U$ alone.  Non-uniformity
requires a non-bijective ingredient in the generator, i.e. an interaction, which
is precisely the role that H1 assigns to the dissipative steps.

\subsection{Measurement as coarse-graining: the Born weights}\label{bridge:born}
\textbf{H2 (fibres).}  The non-bijective steps are coarse only: their fibres
$\pi^{-1}(C)$ --- the sets of ontological states mapped into the same successor
class $C$ --- are exactly the equivalence classes of information that the readable
fields cannot distinguish.

\textbf{H3 (indifference within a fibre).}  The pre-measurement template assigns
equal weight to the ontological states of one fibre: no admissible observer
(Def.~\ref{subsec:observers}) can separate them, and the preparation carries no
information about which member of a fibre will be selected.

\textbf{Theorem 4 (Born rule).}  Under H2 and H3, the probability that a
dissipative step leaves the successor class $C$ readable is
\begin{equation}
p(C)=\sum_{A\in\pi^{-1}(C)}|\lambda_{A}|^{2},
\label{eq:born}
\end{equation}
which is the Born rule; for a uniform template it reduces to the counting measure
$\mu(C)=|\pi^{-1}(C)|/|\mathcal{OS}|$.
\emph{Proof.}  By H2 every ontological state of a fibre has identical readable
fields, so the transition is deterministic on each state and many-to-one onto $C$;
by H3 the ensemble is uniform within a fibre; the frequency of the readable class
$C$ is therefore the counting measure of its fibre, and a general template carries
its linear weight $|\lambda_{A}|^{2}$ by the same measure. \hfill$\blacksquare$

Both hypotheses are the price of the result and should be read as such: H2 is a
property of the rule set, checkable by enumeration at finite $L$; H3 is a statement
about preparation.  What is \emph{not} assumed is the Born rule itself, the
linearity of the evolution, or a collapse postulate.

\subsection{The arrow of time}\label{bridge:arrow}
Under H1-H3 the model needs no collapse postulate beyond the rule: the sieve and
the collapse-and-reissue steps \emph{are} the measurement channel, and their
many-to-one character is the arrow of time of the fabric
(Sect.~\ref{sec-undecidability}).  A template that has not yet been coarse-grained
evolves unitarily by Theorem~2; at a dissipative step it loses the unreadable part
of its weights and continues from the fibre, irreversibly, because the rule is not
injective there.  The model is therefore a completion rather than an
interpretation: the ontological states exist and are all there is, the template is
an observer-level device over them, and the collapse is the fibre map itself.  On a
complete physical reading this is the structure that the superdeterministic
programme of \cite{hossenfelder2020rethinking} describes.

\subsection{No signalling as a corollary}\label{bridge:nosignal}
The readability split that the proposition of Sect.~\ref{subsec:nosignaling-formal}
took as a postulate is H2's fibre structure: the hidden coordination fields are
exactly those that the fibres do not expose.  Causality in the readable sector is
then a corollary of Theorems~2 and~4 --- of the permutation character of the
physical step and the locality of its generator --- rather than a separate
assumption about what observers can see.

\paragraph{Hypotheses, and what remains open even under them.}
\begin{itemize}
\item[(H1)] the fragment boundary (bijective versus many-to-one stages) --- to be
verified by enumeration at finite $L$, not by inspection;
\item[(H2)] the fibres of the dissipative steps equal the unreadable equivalence
classes --- checkable on the reference build;
\item[(H3)] indifference within a fibre --- a preparation statement, tested over
distinct initial templates (the reference dynamics are deterministic, so repeated
runs of one template give no statistics);
\item[(H4)] the stage-by-stage bijection extends to the whole rule set at finite
$L$ (the derivation above is per stage).
\end{itemize}
Even granting H1-H4, three quantitative items remain: the exact fragment boundary
at finite $L$; the decoherence time (how many dissipative steps a template needs
before its off-diagonal weights become unreadable); and whether the fibres
coincide with the physical superselection labels (charge word, sector) or are
finer than them.  Each is a bounded computation on the reference build.

